\chapter{Adaptive Schedules from a Measured Susceptibility}
\label{ch:optimal}

Standard SQA follows a pre-designed $\Gamma(t)$ ramp fixed before the run
begins. Such a schedule cannot react to the system it is driving: it may
rush through the quantum critical region (exciting the system and spoiling
the solution) while wasting steps in regions where nothing interesting
happens. This chapter derives \qv's two alternatives: schedules that
\emph{measure} how critical the simulated system currently is --- through a
susceptibility $\chiB$ --- and pace the annealing trajectory accordingly.
They differ in their order parameter, their regularization, and their
execution model:

\begin{itemize}
\item \textbf{The optimal adaptive $\Jp$ schedule}
(Sections~\ref{sec:roland-cerf}--\ref{sec:optimal-diagnostics}) drives the
Trotter coupling $\Jp$ directly by the local adiabaticity ODE
$d\Jp/dt=\epst\,\chiB^{-\alpha}$, with $\chiB=\operatorname{Var}(B)$ the
susceptibility of the imaginary-time bond sum $B$. Available for \code{sqa}
and \code{sqapt}
(with or without SW cluster moves) via \code{schedule\_type="optimal"}.
\item \textbf{The worldline-magnetization susceptibility schedule}
(Section~\ref{sec:sqa-chi}) instead allocates the fixed step budget over
the annealing parameter $s$ by a floor-regularized weight
$w(s)=\max(\chiB(s),\text{floor})$, where $\chiB$ is now the susceptibility
of the worldline \emph{magnetization} $m$. It is a standalone solver
\emph{method}, \code{method="sqa\_chi"} --- not a \code{schedule\_type} of
\code{sqa} --- with its own parity-parallel update kernel.
\end{itemize}

Section~\ref{sec:sqa-chi-vs-optimal} compares the two directly.

\section{Local adiabaticity: the Roland--Cerf condition}
\label{sec:roland-cerf}

The adiabatic theorem (Section~\ref{sec:qa}) bounds the \emph{total} runtime
by the worst-case gap, Eq.~\eqref{eq:adiabatic}. Roland and
Cerf~\cite{rolandcerf} sharpened this into a \emph{local} condition: rather
than sweeping the control parameter uniformly, demand that the
\emph{instantaneous} transition probability stay uniformly small. For a
control parameter $\lambda(t)$ this reads
\begin{equation}
\Bigl|\frac{d\lambda}{dt}\Bigr|
\;\le\;\eps\,
\frac{\Delta_{\mathrm{gap}}^2(\lambda)}{\bigl|\bra{1}\partial_\lambda\hat H\ket{0}\bigr|}
\;\;\sim\;\;\eps\,\Delta_{\mathrm{gap}}^2(\lambda),
\label{eq:local-adiabatic}
\end{equation}
with $\eps\ll1$ a fixed error budget. The schedule then \emph{slows down
exactly where the gap is small} --- near the quantum critical point --- and
accelerates where the gap is large. For the archetypal Grover problem this
converts the naive $O(1/\Delta_{\min}^2)$ total time into the provably
optimal $O(1/\Delta_{\min})$; generically it redistributes a fixed time
budget in the provably best way.

\subsection{From the gap to a measurable quantity}

Eq.~\eqref{eq:local-adiabatic} is not directly usable: the gap of an
$n$-spin system is unknown. But near a second-order quantum phase
transition, \emph{finite-size scaling} links the gap to correlation
functions that Monte Carlo can measure. With correlation length $\xi$,
dynamical exponent $z$ and anomalous dimension $\eta$:
\begin{equation}
\Delta_{\mathrm{gap}}\sim\xi^{-z},
\qquad\qquad
\chi\;\sim\;\xi^{\,2-\eta},
\end{equation}
where $\chi$ is the susceptibility of the operator driving the transition.
Eliminating $\xi$:
\begin{equation}
\Delta_{\mathrm{gap}}^2 \;\sim\; \chi^{-2z/(2-\eta)} .
\end{equation}
In the Suzuki--Trotter representation the natural control parameter is the
Trotter coupling $\Jp$ (Eq.~\eqref{eq:jperp}) and the natural susceptibility
is that of the quantity $\Jp$ multiplies in the action
Eq.~\eqref{eq:action-cb}: the total imaginary-time \textbf{bond sum}
\begin{keyeq}[Bond sum and bond susceptibility]
\begin{equation}
B=\sum_{t=1}^{M}\sum_{i=1}^{n} s^t_i\,s^{t+1}_i,
\qquad\qquad
\chiB=\operatorname{Var}(B)=\langle B^2\rangle-\langle B\rangle^2 .
\label{eq:chiB}
\end{equation}
\end{keyeq}
$B$ plays the role of the conjugate ``order parameter density'' of the
transverse-field-driven transition; $\chiB$ is cheap to track (see
Section~\ref{sec:chib-est}) and peaks near the quantum transition.
Substituting into Eq.~\eqref{eq:local-adiabatic} and absorbing an extra
$+\tfrac12$ in the exponent that accounts for the statistical width of the
finite-sample variance estimator, the discrete update per adaptive step
becomes

\begin{keyeq}[Adaptive update law]
\begin{equation}
\Delta\Jp \;=\; \epst\cdot\chiB^{-\alpha},
\qquad\qquad
\alpha=\frac{z}{2-\eta}+\frac12 .
\label{eq:update-law}
\end{equation}
\end{keyeq}

For the 1-D quantum Ising universality class ($z=1$, $\eta=\tfrac14$):
$\alpha=\frac{1}{7/4}+\frac12=\frac{4}{7}+\frac12=\frac{15}{14}\approx1.071$
--- the library default (\code{optimal\_alpha}). For the 2-D class use
$\alpha=7/4$; for mean-field (Sherrington--Kirkpatrick-like) problems
$\alpha\approx1$ is a reasonable, if tentative, choice. In practice
$15/14$ is a robust default even when the true universality class is
unknown, because the calibration below makes the schedule's overall pace
independent of moderate errors in $\alpha$.

\begin{physbox}
For this model $\chiB$ is \emph{large near the quantum transition and
collapses toward zero deep in the ordered (classical) phase} --- it peaks
rather than diverging. Through Eq.~\eqref{eq:update-law} the schedule
therefore takes small steps ($\chiB$ large) in the critical window and large
steps ($\chiB$ small) once the system has ordered: a characteristic
\textbf{fast--slow--fast} trajectory in $\Jp$, dwelling exactly where
diabatic errors would otherwise be created.
\end{physbox}

\begin{figure}[ht]
\centering
\begin{tikzpicture}
\begin{axis}[
  width=0.56\textwidth,height=5.4cm,
  xlabel={adaptive step $t$},ylabel={$\Jp$},
  axis lines=left,ticks=none,
  every axis label/.append style={font=\small},clip=false]
\addplot[qnavy,very thick,smooth] coordinates {
  (0,0.1) (0.05,2.5) (0.10,4.8) (0.15,6.8) (0.20,8.4) (0.25,9.6)
  (0.30,10.4) (0.35,10.9) (0.40,11.2) (0.45,11.45) (0.50,11.7)
  (0.55,11.95) (0.60,12.3) (0.65,12.9) (0.70,13.9) (0.75,15.3)
  (0.80,17.0) (0.85,18.9) (0.90,20.8) (0.95,22.6) (1.0,24.0)};
\node[font=\scriptsize\sffamily\color{qgray},anchor=west] at (axis cs:0.47,8.2)
  {dwell near $\Jp^{c}$};
\draw[qamber,densely dashed] (axis cs:0.30,0) -- (axis cs:0.30,22);
\draw[qamber,densely dashed] (axis cs:0.66,0) -- (axis cs:0.66,22);
\node[font=\scriptsize\sffamily\color{qamber!80!black}] at (axis cs:0.48,23.2)
  {critical window};
\end{axis}
\begin{scope}[shift={(8.0,0)}]
\begin{axis}[
  width=0.42\textwidth,height=5.4cm,
  xlabel={$\Jp$},ylabel={$\chiB$},
  axis lines=left,ticks=none,
  every axis label/.append style={font=\small},
  domain=0:1,samples=120]
\addplot[qteal,very thick] {exp(-((x-0.42)^2)/(2*0.018)) + 0.06*exp(-3*x)};
\node[font=\scriptsize\sffamily\color{qgray}] at (axis cs:0.42,1.12)
  {peak at $\Jp^{c}\!\approx\! J_{\mathrm{rms}}$};
\end{axis}
\end{scope}
\end{tikzpicture}
\caption{Schematic of a healthy adaptive run. Left: the realised
$\Jp(t)$ is fast--slow--fast, dwelling in the critical window. Right: the
bond susceptibility $\chiB(\Jp)$ that drives the pacing, peaking near the
quantum transition and collapsing in the ordered phase.}
\label{fig:fastslow}
\end{figure}

\section{Budget calibration: fixing the scale of \texorpdfstring{$\epst$}{eps-tilde}}
\label{sec:calibration}

Eq.~\eqref{eq:update-law} separates cleanly into \emph{shape} and
\emph{scale}: the factor $\chiB^{-\alpha}$ dictates \emph{where} the schedule
is slow or fast, while $\epst$ only sets the overall pace. Choosing
$\epst$ by hand is fragile: since $\chiB$'s magnitude grows with $n$, $M$
and the coupling distribution, a value that works for one instance can make
another instance's schedule either leap across its whole range in a few
steps or --- the more insidious failure --- \emph{freeze}, advancing $\Jp$ so
slowly that the run effectively anneals at constant coupling. Neither
failure announces itself: the run completes and returns an energy.

\qv{} therefore ties $\epst$ to the \emph{step budget}. Treat
Eq.~\eqref{eq:update-law} as an ODE in continuous step-time,
$\frac{dJ}{dt}=\epst\,\chiB(J)^{-\alpha}$, and separate variables:
$dt=\epst^{-1}\chiB(J)^{\alpha}\,dJ$. Requiring the traversal of
$[J_0,J_{\mathrm{end}}]$ to take \emph{exactly} $T=\code{num\_steps}$ steps,
\begin{equation}
T=\frac{1}{\epst}\int_{J_0}^{J_{\mathrm{end}}}\chiB(J)^{\alpha}\,dJ
\end{equation}

\begin{keyeq}[Budget-calibrated adiabaticity scale]
\begin{equation}
\epst \;=\; \frac{1}{T}\int_{J_0}^{J_{\mathrm{end}}}\chiB(J)^{\alpha}\,dJ .
\label{eq:calibrated-eps}
\end{equation}
\end{keyeq}

\begin{warnbox}[Sign of the exponent]
Note the \emph{positive} power $\chiB^{+\alpha}$ inside the integral and the
$1/T$ prefactor. The superficially plausible alternative
$\epst=\Delta\Jp/\sum_m\chiB(J_m)^{-\alpha}$ (negative power, no $1/T$) is
\emph{not} budget-correct when $\chiB$ varies along the range --- it under- or
over-shoots and can re-freeze the schedule. Eq.~\eqref{eq:calibrated-eps}
follows uniquely from the separation of variables above.
\end{warnbox}

\subsection{The pilot pre-pass}

$\chiB(J)$ is not known in advance, so when \code{eps\_tilde}\,$\le0$
(the default \code{optimal\_eps\_tilde=0.0} requests exactly this)
\code{run\_optimal} first runs a short \emph{pilot}: it measures $\chiB$ at
\code{calib\_probes} (default 12) values of $\Jp$ spanning
$[J_0,J_{\mathrm{end}}]$, using \code{calib\_sweeps} (default 10) sweeps per
probe with the \emph{same} sampler the real run will use, on scratch states
so the production run is unbiased. The integral in
Eq.~\eqref{eq:calibrated-eps} is then evaluated by the midpoint rule over
the probe grid.

\subsection{The inverse-CDF trajectory}

One more robustness step: rather than integrating
Eq.~\eqref{eq:update-law} \emph{online} with the noisy within-run $\chiB$
(whose fluctuations can compound --- a run that momentarily orders faster than
the pilot predicted sees a small $\chiB$, takes a huge step, and skips the
rest of the range), the entire trajectory is \emph{precomputed} from the
pilot profile. Define the cumulative ``adiabatic cost''
\begin{keyeq}[Inverse-CDF placement of steps]
\begin{equation}
G(J)=\int_{J_0}^{J}\chiB(J')^{\alpha}\,dJ' ,
\qquad\qquad
J^{(t)}=G^{-1}\!\Bigl(\frac{t}{T-1}\,G(J_{\mathrm{end}})\Bigr),
\quad t=0,\dots,T-1 .
\label{eq:inverse-cdf}
\end{equation}
\end{keyeq}
Each step is placed at an equal increment of adiabatic cost. By
construction --- independent of the shape or scale of $\chiB$ --- the schedule
(a)~consumes exactly $T$ steps, (b)~reaches $J_{\mathrm{end}}$, and
(c)~concentrates its steps where $\chiB$ is large. The online $\chiB$ is
still measured at every step and recorded in \code{chi\_B\_trace} as a
diagnostic (and, in the debug CSV, as \code{chi\_B\_online}).

A \emph{positive} \code{eps\_tilde} bypasses calibration entirely and drives
the legacy online update $\Jp\mathrel{+}=\epst\,\chiB^{-\alpha}$ with that
fixed scale --- retained for reproducibility and for controlled experiments,
not recommended for production.

\subsection{Fairness of the step budget}

The adaptive run \emph{always} executes all \code{num\_steps} steps; if
$\Jp$ saturates at $J_{\mathrm{end}}$ early it keeps sweeping at the
endpoint rather than stopping. Total sweep effort is therefore identical to
a standard-schedule run of the same length --- any quality difference is
attributable to the \emph{placement} of the steps, not to extra compute.

\section{Estimating \texorpdfstring{$\chiB$}{chi\_B} without overhead}
\label{sec:chib-est}

\subsection{Incremental bond tracking}

Recomputing $B$ (Eq.~\eqref{eq:chiB}) from scratch costs $O(nM)$; instead,
every move type updates it incrementally in $O(1)$ per flip:

\begin{center}
\small
\begin{tabular}{@{}lp{9.1cm}@{}}
\toprule
\textbf{Move} & \textbf{Bond-sum change $\Delta B$}\\
\midrule
Metropolis flip of $s^t_i$ &
$\Delta B=-2\,s^t_i\bigl(s^{t-1}_i+s^{t+1}_i\bigr)$ evaluated with the
\emph{old} spin value --- the two touched time-bonds reverse.\\
SW cluster flip &
$\Delta B=-2\sum_{\text{boundary bonds}} s\,s'$ --- only the bonds joining
the cluster to unflipped spins change; interior bonds are invariant.\\
Worldline flip &
$\Delta B=0$ --- every bond of the worldline is a product of two flipped
spins.\\
\bottomrule
\end{tabular}
\end{center}

After each full sweep the current $B$ is appended to the step's sample
buffer, so $\chiB$ estimation adds essentially zero cost to the run.

\subsection{Pooled variance across sweeps and replicas}

Each adaptive step yields $R\times S$ samples of $B$ ($R$ replicas,
$S$ sweeps within the step), pooled into one estimate
$\chiB=\widehat{\operatorname{Var}}(B)$. The pooling is statistically
well-behaved in both regimes:
\begin{itemize}
\item \emph{Near criticality}, the autocorrelation time
$\tau_{\mathrm{int}}\gg S$, so each replica's $S$ samples are nearly
identical and the pooled variance collapses to the \emph{cross-replica}
variance --- exactly the honest estimator for correlated chains.
\item \emph{Away from criticality}, $\tau_{\mathrm{int}}\ll S$ and the
pooled estimate enjoys $\sim RS$ effectively independent samples.
\end{itemize}
This is the reason \code{replicas}\,$\geq4$ is recommended whenever
\code{schedule\_type="optimal"} is used.

\subsection{Single-replica guard and the \texorpdfstring{$\chiB$}{chi-B} floor}

Two safeguards prevent pathological steps:
\begin{description}
\item[$R=1$ guard.] With a single replica near criticality all within-step
samples coincide, the naive variance $\to0$, and
Eq.~\eqref{eq:update-law} would take its \emph{maximum} step precisely where
it must be smallest. The implementation therefore uses
$\chiB=\max\bigl(\widehat{\operatorname{Var}}(B),\;
nM\bigl[1-(B/nM)^2\bigr]\bigr)$, the second term being the
single-bond variance extrapolated over the lattice --- a strictly positive
lower bound near criticality.
\item[Floor.] Deep in the ordered phase $\chiB\to0$ and
$\chiB^{-\alpha}$ blows up; a floor
$\chiB\ge10^{-4}\max_J\chiB^{\text{(pilot)}}$ caps the step size. The debug
CSV records where the floor activates (\code{floor\_hit}); in a healthy run
this is only the deeply ordered tail, never the critical window.
\end{description}

\section{Endpoints, temperature, and the two-phase SQA protocol}
\label{sec:optimal-params}

\subsection{Choosing $J_0$ and $J_{\mathrm{end}}$}

All endpoints derive from the problem's own coupling scale. Let
\begin{equation}
J_{\mathrm{rms}}=\sqrt{\operatorname{mean}\bigl(J_{ij}^2\bigr)}
\quad\text{over the non-zero couplings}
\end{equation}
(computable via \code{j\_rms\_from\_problem(problem)}). The quantum critical
coupling scales as $\Jp^{c}\sim J_{\mathrm{rms}}$, so:
\begin{itemize}
\item \code{j\_perp\_start} --- from the tuned standard schedule's
$(\beta,\Gamma_{\text{start}})$ through Eq.~\eqref{eq:jperp}; small
(quantum end).
\item \code{j\_perp\_end} --- default $\max(J_0,\;5\,J_{\mathrm{rms}})$,
comfortably past the transition (classical end).
\item $\beta$ --- for SQA, $\beta=\max(4/J_{\mathrm{rms}},\,1)$, i.e.\ cold
enough that $\beta J_{\mathrm{rms}}\geq4$: thermal fluctuations must not
wash out the order the $\Jp$ ramp is trying to establish.
\end{itemize}
The helper returns all three at once:
\begin{pycode}
from qanneal import optimal_j_perp_params, j_rms_from_problem
beta, jp_start, jp_end = optimal_j_perp_params(problem, trotter_slices=32)
jr = j_rms_from_problem(problem)
\end{pycode}

\begin{warnbox}[When can the adaptive schedule help at all?]
The schedule needs a quantum range to act on: it helps when
$J_{\mathrm{rms}}\gtrsim\code{j\_perp\_start}$ (in the auto-tuned setting,
roughly $J_{\mathrm{rms}}\gtrsim3$), so that the run genuinely starts on the
disordered side of the transition. If $J_{\mathrm{rms}}$ is far below
$J_0$ the system is already effectively classical at step one,
$\chiB\approx0$ everywhere, and no schedule --- adaptive or not --- has
anything to optimise. Check a new problem class with
\code{scripts/sanity\_schedule.py}, which classifies instances as
\code{PASS} or \code{DEGENERATE} before you spend compute.
\end{warnbox}

\begin{warnbox}[Always pass an explicit \code{j\_perp\_end} in scripted runs]
The C++ \code{run\_optimal} accepts a sentinel \code{j\_perp\_end<=0}
that resolves to the coldest replica's Trotter coupling. If that resolves
\emph{below} \code{j\_perp\_start} the schedule has no range; the code then
installs a large fallback ceiling and prints a \code{cerr} warning. Relying
on the sentinel makes runs silently sensitive to ladder details --- pass an
explicit \code{j\_perp\_end} (e.g.\ $5J_{\mathrm{rms}}$) in anything
scripted.
\end{warnbox}

\subsection{Two-phase protocol for SQA}

Plain SQA under the adaptive schedule holds a \emph{single} fixed $\beta$
--- it has no thermal annealing at all, which measurably hurts it. The SQA
\code{run\_optimal} therefore runs two phases, split by
\code{beta\_ramp\_fraction} (default $0.3$):
\begin{enumerate}
\item \textbf{Thermal ramp} --- for the first
$\lceil0.3\cdot\code{num\_steps}\rceil$ steps, $\beta$ ramps from
\code{beta\_ramp\_start} (default $0.1$) up to the cold target
$\beta=\max(4/J_{\mathrm{rms}},1)$ at \emph{fixed} $\Jp=J_0$: a conventional
thermal anneal that settles the slices into the disordered quantum state at
the right temperature.
\item \textbf{Adaptive ramp} --- $\beta$ then stays fixed while $\Jp$
follows the calibrated inverse-CDF trajectory
Eq.~\eqref{eq:inverse-cdf} over the remaining steps.
\end{enumerate}
SQAPT does \emph{not} use phase 1: its $(\beta,\Gamma)$ ladder plus replica
swaps already supply thermal mobility.

\subsection{SQAPT under a shared $\Jp$}
\label{sec:sqapt-optimal}

In the SQAPT variant, \emph{all} rungs share the single evolving $\Jp$;
each keeps its own $(\beta_r,\Gamma_r)$ from the ladder. Two consequences,
both derived earlier:
\begin{itemize}
\item In the swap exponent Eq.~\eqref{eq:swap-split} the quantum term
carries the factor $\Jp^{(r)}-\Jp^{(r+1)}=0$: \textbf{swaps reduce to the
purely classical criterion} on slice-summed energies.
\item $\chiB$ is estimated by pooling $B$-samples across all rungs, exactly
as in Section~\ref{sec:chib-est}.
\end{itemize}

\section{Usage}

\subsection{High-level}

\begin{pycode}
from qanneal import solve

result = solve(
    problem,
    method="sqa",              # or "sqapt"
    schedule_type="optimal",
    reads=15,
    replicas=4,                # >=4 for a reliable chi_B estimate
    optimal_eps_tilde=0.0,     # 0.0 = per-instance budget calibration
    optimal_num_steps=150,
    # Optional overrides (auto-computed from the problem scale if omitted):
    # optimal_j_perp_end=25.0,         # default max(j_perp_start, 5*j_rms)
    # optimal_beta_ramp_fraction=0.3,  # SQA two-phase split
    # optimal_alpha=15/14,
    # optimal_debug_csv="schedule.csv" # per-step J_perp / chi_B diagnostics
)
print(result.best_energy, result.best_sample)
\end{pycode}

For SQAPT with SW cluster moves:

\begin{pycode}
result = solve(
    problem,
    method="sqapt",
    schedule_type="optimal",
    replicas=8, reads=15,
    trotter_slices=32,
    worldline_sweeps=3,
    cluster_sweeps=1,          # >0 = SQAPT+SW
    swap_interval=1,
    optimal_eps_tilde=0.0,
    optimal_num_steps=150,
)
\end{pycode}

\subsection{Low-level, with full diagnostics}

\begin{pycode}
from qanneal import SQAAnnealer, SQASchedule

dummy = SQASchedule.from_vectors([1.0], [1.0])  # construction only
ann = SQAAnnealer(ising, dummy, trotter_slices=32, replicas=4)

result = ann.run_optimal(
    beta=1.0,                    # phase-2 (cold) inverse temperature
    j_perp_start=0.1, j_perp_end=25.0,
    eps_tilde=0.0,               # <=0 -> budget calibration
    alpha=15/14, num_steps=150, sweeps_per_step=20,
    worldline_sweeps=3, cluster_sweeps=0,
    calib_probes=12, calib_sweeps=10,   # pilot grid
    beta_ramp_fraction=0.3,             # two-phase thermal ramp
)
result.j_perp_trace          # realised J_perp(t)
result.chi_B_trace           # online chi_B(t)
result.calibrated_eps_tilde  # the eps from Eq. (calibration)
result.final_j_perp          # J_perp actually reached
result.best_state, result.best_energy
\end{pycode}

The SQAPT twin is \code{SQAParallelTemperingAnnealer.run\_optimal}
(full signature in Chapter~\ref{ch:api}), with identical calibration
semantics and no two-phase ramp.

\section{Diagnosing a run}
\label{sec:optimal-diagnostics}

Plot the three traces after any adaptive run:

\begin{pycode}
import numpy as np, matplotlib.pyplot as plt

t  = np.arange(len(result.j_perp_trace))
dj = np.diff(result.j_perp_trace)

fig, axs = plt.subplots(3, 1, figsize=(8, 9), sharex=True)
axs[0].plot(t, result.j_perp_trace);  axs[0].set_ylabel("J_perp")
axs[1].plot(t[:-1], dj);              axs[1].set_ylabel("step size dJ_perp")
axs[2].plot(t, result.energy_trace);  axs[2].set_ylabel("energy")
axs[2].set_xlabel("adaptive step")
plt.suptitle("Optimal schedule diagnostics"); plt.tight_layout(); plt.show()
\end{pycode}

\textbf{A healthy run} shows (Fig.~\ref{fig:fastslow}): fast early movement,
a pronounced dwell (small $\Delta\Jp$) in the critical window near
$\Jp\approx J_{\mathrm{rms}}$, fast movement again past the transition, the
endpoint reached within budget, and the steepest energy descent during the
dwell.

\textbf{Warning signs:}
\begin{itemize}
\item \emph{Flat or linear $\Jp(t)$ with no dwell} --- the schedule is
frozen or degenerate. Check that $J_{\mathrm{end}}>J_0$ (look for the
\code{cerr} fallback warning) and that the problem is in the favourable
regime ($J_{\mathrm{rms}}\gtrsim J_0$).
\item \emph{Constant $\Delta\Jp$} --- $\chiB$ estimate too noisy to shape
the trajectory; raise \code{replicas} and/or \code{sweeps\_per\_step}.
\item \emph{\code{final\_j\_perp} far below \code{resolved\_j\_perp\_end}}
on a non-degenerate instance --- calibration failed; inspect the debug CSV
(columns \code{phase}, \code{step\_index}, \code{j\_perp}, \code{chi\_B},
\code{delta\_j}, \code{eps\_tilde}, \code{alpha}, \code{floor\_hit},
\code{chi\_B\_online}; \code{phase=calib} rows are pilot probes).
The helper \code{scripts/}\allowbreak\code{plot\_schedule\_trajectory.py}
renders it.
\end{itemize}

\begin{notebox}[Pre-flight habit]
Before any long or large-$n$ run, spend a minute on a small instance of the
same problem class: run with \code{optimal\_debug\_csv} set, plot the
trajectory, and confirm the fast--slow--fast shape.
\code{python scripts/sanity\_schedule.py --n 40 --steps 150 --plot}
automates exactly this check and prints \code{PASS} / \code{DEGENERATE} /
\code{FAIL} per instance.
\end{notebox}

\section{Attribution: schedule effect vs.\ cluster moves}
\label{sec:attribution}

The adaptive \emph{schedule} and SW \emph{cluster sweeps} are independent
mechanisms that both help on hard instances. When evaluating the schedule on
your own problems, keep them separate or you will misattribute gains:

\begin{center}
\small
\begin{tabular}{@{}lll@{}}
\toprule
\textbf{Comparison} & \textbf{\code{cluster\_sweeps}} & \textbf{Isolates}\\
\midrule
\code{sqapt}-std vs \code{sqapt}-opt & $0$ & the schedule effect alone\\
\code{sqapt}+SW-std vs \code{sqapt}+SW-opt & $>0$ & schedule $\times$ cluster interaction\\
\bottomrule
\end{tabular}
\end{center}

Only credit ``the optimal schedule'' if the first comparison reproducibly
favours it; if only the +SW pair improves, the gain belongs to the
interaction and should be framed as such.

\section{The worldline-magnetization susceptibility schedule
(\texorpdfstring{\code{sqa\_chi}}{sqa\_chi})}
\label{sec:sqa-chi}

The schedule above measures $\chiB$ from the imaginary-time bond sum $B$
and drives $\Jp$ online through \code{SQAAnnealer}. \qv{} 1.0.0 adds a
second, \emph{standalone} method built around a different order parameter
and a different execution model: \code{sqa\_chi} is not a
\code{schedule\_type} flag on \code{sqa} but its own solver class,
\code{SQAChiAnnealer}, generalized from a worldline-QMC validation script
into \qv's \code{Backend} abstraction (dense \emph{or} sparse Ising).

\subsection{Order parameter: the worldline magnetization}
\label{sec:sqa-chi-order-param}

Instead of the bond sum $B$ (Eq.~\eqref{eq:chiB}), \code{sqa\_chi} measures
the mean spin over the \emph{entire} worldline lattice of a replica ($n$
spins $\times\,M$ Trotter slices):
\begin{keyeq}[Worldline magnetization and its susceptibility]
\begin{equation}
m=\frac{1}{nM}\sum_{i=1}^{n}\sum_{k=1}^{M}\sigma_i^k,
\qquad\qquad
\chiB = nM\bigl(\langle m^2\rangle-\langle|m|\rangle^2\bigr),
\label{eq:chib-magnetization}
\end{equation}
\end{keyeq}
the ordinary classical-Ising susceptibility of the $(d{+}1)$-dimensional
Suzuki--Trotter lattice. Like $\operatorname{Var}(B)$, it peaks at the
quantum critical point --- both are measuring the growth of critical
fluctuations in the \emph{same} classical system, just through different
observables of it. $\langle\cdot\rangle$ pools every post-burn sweep sample
across all replicas, mirroring Section~\ref{sec:chib-est}'s pooling
(replicas $\ge4$ recommended for the same statistical reasons).

\subsection{The annealing parameter \texorpdfstring{$s$}{s}}

\code{sqa\_chi} works in the annealing parameter $s\in[0,1]$ of the
symmetric path $A(s)=A_0(1-s)$, $B(s)=s$ (driver and problem Hamiltonian
weights), rather than in $\Gamma$ directly:
\begin{keyeq}[Annealing parameter $\leftrightarrow$ transverse field]
\begin{equation}
s=\frac{A_0}{A_0+\Gamma},
\qquad\qquad
\Gamma=\frac{A_0(1-s)}{s},
\label{eq:sqa-chi-s}
\end{equation}
\end{keyeq}
so $s=0$ is the fully quantum end ($\Gamma\to\infty$) and $s=1$ the fully
classical end ($\Gamma=0$); the driver scale $A_0$ is \code{chi\_driver\_A0}
if positive, otherwise \code{gamma\_start}. This is a property of the
annealing \emph{path}, not of the $\chiB$ estimator --- the same relation
underlies every $s$-parametrized construction in \qv.

\subsection{Stage 1: the pilot \texorpdfstring{$\chiB(s)$}{chi\_B(s)} scan}

At \code{scan\_points} (default 16) uniform values of $s$ between
$s(\gamma_{\mathrm{start}})$ and $s(\gamma_{\mathrm{end}})$, the pilot runs
\code{scan\_burn}$+$\code{scan\_sweeps} (default $10+30$) parity-parallel
checkerboard sweeps (Section~\ref{sec:sqa-chi-kernel}) at the corresponding
fixed $\Jp$.

\begin{warnbox}[Fresh worldline per grid point --- deliberately not carried]
Unlike \code{run\_optimal}'s calibration pilot and the retired surrogate
schedule (both of which carry the worldline state between grid points so
the scan itself gently anneals), \code{sqa\_chi} starts every grid point
from a \emph{fresh, independently randomized} worldline. Each $\chiB(s)$
estimate is therefore a statistically independent measurement --- more
robust against carried-over bias between probes, at the cost of
re-equilibrating (\code{scan\_burn} sweeps) at every point.
\end{warnbox}

The scan profile (\code{scan\_s}, \code{scan\_gamma}, \code{scan\_chi\_B})
and the quantum critical point estimate --- the $\chiB$ peak, recorded as
(\code{s\_star}, \code{gamma\_star}, \code{j\_perp\_star}) --- are returned
as diagnostics.

\subsection{Stage 2: floor-regularized time allocation}

The pilot profile is turned into a floor-regularized weight and its
cumulative integral:
\begin{keyeq}[Floored weight and its cumulative integral]
\begin{gather}
w(s)=\max\bigl(\chiB(s),\,\text{floor}\bigr),
\qquad
\text{floor}=\max\Bigl(\code{chi\_floor\_fraction}\cdot\max_s\chiB(s),\,
10^{-12}\Bigr),
\label{eq:sqa-chi-weight}\\[4pt]
\tau(s)=\frac{\displaystyle\int_{s_{\mathrm{lo}}}^{s}w(s')\,ds'}
             {\displaystyle\int_{s_{\mathrm{lo}}}^{s_{\mathrm{hi}}}w(s')\,ds'} .
\end{gather}
\end{keyeq}
$\chiB(s)$ linearly interpolated between scan points (clamped to the
boundary value outside $[s_{\mathrm{lo}},s_{\mathrm{hi}}]$). Each of the
remaining (post-ramp) adaptive steps $k=0,\dots,n_2-1$ is placed at
$s_k=\tau^{-1}\!\bigl(k/(n_2-1)\bigr)$, evaluated on a fixed 4096-point grid
by trapezoidal integration and linear inversion, then mapped to
$\gamma_k=A_0(1-s_k)/s_k$ (clamped into
$[\gamma_{\mathrm{end}},\gamma_{\mathrm{start}}]$).

\begin{notebox}[A floor, not an additive regularizer]
The optimal schedule's calibration floor and the retired surrogate
schedule both use an \emph{additive} term ($\chiB+\chi_0$) to keep the
weight strictly positive. \code{sqa\_chi} instead \emph{clamps}
$\chiB$ from below --- a multiplicative floor, matching the reference
validation script's regularization idea (\code{np.maximum(chi\_B, eps)})
but generalized: the script's fixed absolute $\epsilon=10^{-8}$ is not
scale-invariant across problem sizes ($\chiB\sim O(nM)$), so \qv{} scales
the floor to a fraction of the scan's own peak instead.
\end{notebox}

\begin{notebox}[Degenerate profile]
If the pilot scan is flat or non-finite (no measurable $\chiB$ signal),
\code{floor} falls back to $1.0$ and a \code{cerr} warning is printed;
$w(s)$ is then constant and the schedule degrades gracefully to a
linear-in-$s$ ramp rather than failing.
\end{notebox}

\subsection{The parity-parallel checkerboard kernel}
\label{sec:sqa-chi-kernel}

Every sweep in \code{sqa\_chi} --- pilot \emph{and} production --- uses a
kernel that differs from every other engine in \qv{}: it parallelizes
\emph{within} a single replica's worldline, not just across replicas.

\begin{description}[style=nextline]
\item[The observation.] Trotter slices of the same parity (even/odd index)
are never directly coupled: the Trotter term only links slice $k$ to
$k\pm1$, which always belong to the \emph{other} parity. So while one
parity's slices are held fixed, every $(\text{replica},\text{slice})$ cell
of the other parity can be updated completely independently.
\item[The update.] Each cell runs a full, ordinary sequential single-spin
Metropolis sweep over its own $n$ spins, using the \code{Backend}'s cached
local fields (Section~\ref{sec:chib-est}'s $O(1)$/$O(\text{degree})$
incremental updates) --- exactly the update rule \code{SQAAnnealer} uses,
just scoped to one slice at a time.
\item[The parallelism.] All cells of one parity are dispatched to a single
OpenMP \code{\#pragma omp parallel for}, then the other parity, then the
lattice's pooled magnetization is read off. Two parity passes make one
full sweep.
\end{description}

\begin{warnbox}[Why not update a whole slice's spins simultaneously?]
A natural-looking shortcut --- flip every spin of a slice at once, using
the \emph{previous} sweep's neighbour values for the in-plane coupling ---
is what the reference validation script does. It is only exact detailed
balance when the spatial coupling graph is \emph{bipartite}; for a general
(dense or frustrated) Ising problem it silently biases the sampler. The
per-slice \emph{sequential} update above is exact for any coupling graph,
while still parallelizing across slices --- the price is that spins
\emph{within} one slice are still updated one at a time, not vectorised.
\end{warnbox}

Because the parallel axis is $(\text{replica}\times\text{slice})$ rather
than replica alone, \code{sqa\_chi} scales across cores even at
\code{replicas=1} — unlike every other SQA-family engine in \qv{}, whose
OpenMP parallelism is over replicas only and is therefore inert for a
single-replica run.

\subsection{Two-phase protocol and the main anneal}

As in Section~\ref{sec:optimal-params}, the first
$\lceil\code{beta\_ramp\_fraction}\cdot\code{num\_steps}\rceil$ steps ramp
$\beta$ from \code{beta\_ramp\_start} to the target $\beta$ at fixed
$\gamma=\gamma_{\mathrm{start}}$ (\code{beta\_ramp\_fraction}$\le0$ disables
this phase); the remaining steps hold $\beta$ fixed and follow the
$\gamma_k$ schedule from Stage 2. The resolved
(\code{beta\_schedule}, \code{gamma\_schedule}) is run on a \emph{fresh}
worldline (independent of the pilot scan's state, so the production run is
unbiased) through the same parity-parallel checkerboard kernel.

\subsection{Usage}

\begin{pycode}
from qanneal import solve

result = solve(
    problem,
    method="sqa_chi",                  # standalone method, not schedule_type="..."
    reads=15,
    replicas=4,                        # >=4 for a reliable chi_B estimate
    optimal_num_steps=150,             # shared step budget with schedule_type="optimal"
    chi_floor_fraction=1e-6,
    # Optional overrides (auto from the resolved standard schedule if omitted):
    # chi_gamma_start=5.0,             # default: max gamma of the tuned schedule
    # chi_gamma_end=0.01,              # default: min gamma of the tuned schedule
    # chi_scan_points=16, chi_scan_sweeps=30, chi_scan_burn=10,
    # chi_driver_A0=0.0,               # <=0 resolves to chi_gamma_start
    # optimal_beta_ramp_fraction=0.3,  # shared two-phase split
    # optimal_debug_csv="sqa_chi_schedule.csv",
)
print(result.best_energy, result.best_sample)
\end{pycode}

Low-level, with full diagnostics:

\begin{pycode}
from qanneal import SQAChiAnnealer

ann = SQAChiAnnealer(ising, trotter_slices=32, replicas=4)   # no schedule argument

result = ann.run_chi(
    beta=1.0, gamma_start=5.0, gamma_end=0.01,
    num_steps=150, sweeps_per_step=20,
    scan_points=16, scan_sweeps=30, scan_burn=10,
    chi_floor_fraction=1e-6, driver_A0=0.0,
    beta_ramp_fraction=0.3, beta_ramp_start=0.1,
)
result.scan_s, result.scan_gamma, result.scan_chi_B     # pilot chi_B(s) profile
result.s_star, result.gamma_star, result.j_perp_star    # QCP estimate
result.beta_schedule, result.gamma_schedule, result.s_schedule  # resolved trajectory
result.best_state, result.best_energy
\end{pycode}

\subsection{Diagnosing a run}

The same three-panel plot as Section~\ref{sec:optimal-diagnostics} applies
to \code{gamma\_schedule} / \code{s\_schedule} / \code{energy\_trace}; also
plot \code{scan\_chi\_B} against \code{scan\_s} to confirm the pilot found a
clear peak (a flat scan is the \code{sqa\_chi} analogue of a frozen
adaptive schedule --- see the degenerate-profile fallback above).

When \code{debug\_csv\_path} is set, columns are \code{phase}
(\code{scan}, \code{beta\_ramp}, or \code{run}), \code{index}, \code{s},
\code{gamma}, \code{j\_perp}, \code{chi\_B}, and \code{beta}; only
\code{scan} rows carry a measured \code{chi\_B} --- \code{beta\_ramp}/
\code{run} rows report the resolved schedule that \code{chi\_B} produced,
not an online remeasurement.

\subsection{\code{sqa\_chi} vs.\ the optimal \texorpdfstring{$\Jp$}{J-perp} schedule}
\label{sec:sqa-chi-vs-optimal}

\begin{center}
\small
\begin{tabular}{@{}p{2.6cm}p{6.0cm}p{6.0cm}@{}}
\toprule
& \textbf{Optimal $\Jp$ schedule} & \textbf{\code{sqa\_chi}}\\
\midrule
Exposed as & \code{schedule\_type="optimal"} on \code{sqa}/\code{sqapt} & standalone \code{method="sqa\_chi"}\\
Order parameter & bond sum $B=\sum_{k,i}\sigma_i^k\sigma_i^{k+1}$ & worldline magnetization $m=\text{mean}(\sigma)$\\
$\chiB$ formula & $\operatorname{Var}(B)$ & $nM(\langle m^2\rangle-\langle|m|\rangle^2)$\\
Control variable & $\Jp$ directly & $s$ (hence $\Gamma$)\\
Regularizer & multiplicative floor on $\chiB^{-\alpha}$ & multiplicative floor on $\chiB$ itself\\
Universality exponent & $\alpha$ required (default $15/14$) & none\\
Pilot-scan state & carried between grid points & fresh worldline per grid point\\
Update kernel & sequential per-slice, parallel over replicas & parity-parallel checkerboard over slices \emph{and} replicas\\
Parallel at \code{replicas=1}? & no & yes\\
Backends & dense or sparse & dense or sparse\\
Methods & \code{sqa}, \code{sqapt} (+SW) & \code{sqa\_chi} only\\
\bottomrule
\end{tabular}
\end{center}

Both precompute their trajectory from a pilot $\chiB$ profile rather than
integrating online, so neither is vulnerable to the compounding-noise
failure mode of Section~\ref{sec:calibration}. Treat the choice between
them as an empirical one: neither dominates the other on every instance,
and Section~\ref{sec:attribution} applies here as well --- isolate the
schedule effect from cluster-move effects (where applicable) before
crediting either one.
